\section{Variantes de la red Q profunda y comparación de entrenamiento}

\subsection{Análisis de las variantes principales}

\begin{table}[H]
\centering
\begin{tabular}{lll}
\toprule
Variante & Mecanismo central & Efecto \\
\midrule
Double DQN & La red actual elige la acción y la red objetivo la evalúa & Reduce la overestimation y mejora la estabilidad \\
Dueling DQN & Descompone $Q=V+A$ & Aprende más rápido el valor del estado; la ventaja de las acciones puede compartir parámetros \\
Prioritized Replay & Muestreo ponderado según el TD error & Recupera pronto las transiciones poco frecuentes pero decisivas \\
Rainbow & Combina varias técnicas de mejora (NoisyNet, Multi-step, Distributional) & Mayor cobertura, pero el ajuste de hiperparámetros es más complejo \\
\bottomrule
\end{tabular}
\caption{Variantes de DQN y su efecto en la ingeniería}
\end{table}

\begin{knowledgebox}{Escenarios adecuados para la combinación de varias técnicas}
No todos los escenarios necesitan el conjunto completo de mejoras de Rainbow. En entornos de baja capacidad, Double DQN combinado con Dueling basta para contener la sobreestimación. En tareas de Atari con abundantes recursos de cómputo o en la conducción simulada, los retornos de varios pasos de Rainbow y NoisyNet aportan una mayor capacidad de generalización.
\end{knowledgebox}

\subsection{Comparación de costos de entrenamiento e inferencia}

\begin{importantbox}{Modelo de costos: entrenamiento frente a inferencia}
El entrenamiento se concentra en el muestreo y la optimización: $Cost_{train} \approx N_{\text{steps}} \times (^{\text{forward}}C + ^{\text{backward}}C)$; la inferencia, en cambio, busca un equilibrio entre rendimiento y latencia: $Cost_{infer} \approx C_{\text{forward}} + C_{\text{action selection}}$. Si se lleva DQN a dispositivos embebidos, conviene considerar la cuantización o la poda para reducir el cómputo.
\end{importantbox}

\begin{knowledgebox}{Desglose de la pila de inferencia}
Un inference pipeline suele incluir: la percepción del entorno (state encoding), la pasada hacia adelante de la red Q, la selección de la acción con $\arg\max$, el filtrado de acciones (por ejemplo, con restricciones previas) y las verificaciones de seguridad en la interfaz con el simulador. En la coordinación multi-agent también hay que controlar los retrasos de comunicación y de sincronización.
\end{knowledgebox}

\begin{warningbox}{Cuantización excesiva en la inferencia de baja latencia}
Para optimizar la latency de inferencia, los equipos de ingeniería suelen cuantizar o podar la red; sin embargo, si no se preserva la ruta crítica (por ejemplo, la rama de advantage de la última capa), el orden de los valores Q puede alterarse. Se recomienda mantener dos canales: la versión cuantizada para la inferencia y un respaldo en full precision para la periodic calibration.
\end{warningbox}

\subsection{Resumen de la sección}

Esta sección, mediante la comparación de las variantes de DQN y el análisis del modelo de costos, señala cómo elegir las técnicas de mejora y la combinación de operadores adecuadas según las distintas restricciones de recursos.
